\section{Prime and Maximal Ideals}

\exercise{27.1}{Prime and maximal ideals in a six-element ring}
Every ideal is $\gen d$ for $d\mid6$, and $\Z_6/\gen d\cong\Z_d$. A proper ideal is prime exactly when its quotient is an integral domain, and maximal exactly when its quotient is a field. Among the divisors $1,2,3,6$, exactly $2,3$ are prime. Thus both lists are $\gen2=\{0,2,4\}$ and $\gen3=\{0,3\}$.

\exercise{27.2}{Prime and maximal ideals in a twelve-element ring}
For $d\mid12$, reduction modulo $d$ gives $\Z_{12}/\gen d\cong\Z_d$. The quotient is a domain, equivalently a field in this finite case, exactly when $d=2$ or $3$. Hence the prime ideals and the maximal ideals are both $\gen2=\{0,2,4,6,8,10\}$ and $\gen3=\{0,3,6,9\}$.

\exercise{27.3}{A product of two two-element fields}
In a product $A\times B$ with identities, every ideal is $I\times J$: multiply $(a,b)$ by $(1,0)$ and $(0,1)$ to separate the coordinates. A domain quotient cannot retain both nonzero factors, since $(1,0)(0,1)=0$. Thus the prime and maximal ideals of $\Z_2\times\Z_2$ are exactly $\{0\}\times\Z_2$ and $\Z_2\times\{0\}$. Each quotient is $\Z_2$.

\exercise{27.4}{A product involving a four-element ring}
The same product argument shows that a prime ideal must be all of one factor times a prime ideal in the other. In $\Z_2$ the only prime ideal is zero, and in $\Z_4$ it is $\gen2=\{0,2\}$. Therefore the only prime ideals are
\[\{0\}\times\Z_4,\qquad\Z_2\times\{0,2\}.\]
Both quotients are $\Z_2$, so these are also exactly the maximal ideals.

\exercise{27.5}{Quadratic quotients over the three-element field}
A quotient by a nonconstant polynomial is a field exactly when that polynomial is irreducible. A quadratic over a field is irreducible exactly when it has no root. In $\Z_3$, the possible squares are $0,1$. The equation $x^2=-c$ has no solution exactly when $-c=2$, or $c=1$. Thus only $c=1$ works.

\exercise{27.6}{A cubic family with fixed quadratic term}
A cubic over a field is reducible exactly when it has a root. At $x=0,1,2$, the values of $x^3+x^2$ modulo three are $0,2,0$. Thus $x^3+x^2+c$ has a root if $c=0$ or $1$, but for $c=2$ its values are $2,1,2$, all nonzero. The required value is $c=2$.

\exercise{27.7}{A cubic family with a variable coefficient}
Evaluate $x^3+cx^2+1$ at $0,1,2$. For $c=0$ the value at $2$ is zero; for $c=1$ the value at $1$ is zero; for $c=2$ the three values are $1,1,2$, all nonzero. Thus exactly $c=2$ makes the cubic irreducible and the quotient a field.

\exercise{27.8}{Quadratics with a linear term over five elements}
Since the characteristic is not two, $x^2+x+c$ has a root exactly when the discriminant $1-4c$ is a square. The squares in $\Z_5$ are $0,1,4$. For $c=0,1,2,3,4$, the discriminants are respectively $1,2,3,4,0$. Thus the nonsquare cases are $c=1,2$, and precisely those quotients are fields.

\exercise{27.9}{Quadratics with constant term one}
For $x^2+cx+1$, the discriminant is $c^2-4$. As $c$ runs through $0,1,2,3,4$ in $\Z_5$, the discriminants are $1,2,0,0,2$. Nonsquares occur for $c=1,4$, so exactly those values give fields.

\exercise{27.10}{Correcting maximal ideal}
A maximal ideal $M$ is a \emph{proper} ideal of $R$ such that the only ideals containing it are $M$ and $R$. Merely saying that no other ideal contains it would make $R$ the only candidate, because $R$ contains every ideal.

\exercise{27.11}{Correcting prime ideal}
In a commutative ring, a prime ideal $P$ is a proper ideal such that $ab\in P$ implies $a\in P$ or $b\in P$ for all $a,b\in R$. It need not be generated by a prime number; for example $(x)\subset\Q[x]$ is prime because its quotient is $\Q$.

\exercise{27.12}{Prime field}
The definition is correct: a prime field has no proper subfields. Equivalently it equals the subfield generated by its unity. Up to isomorphism these fields are $\Q$ and $\Z_p$ for primes $p$.

\exercise{27.13}{Principal ideal}
The definition is correct. In a commutative ring with unity, the smallest ideal containing $a$ is $(a)=aR=\{ar:r\in R\}$. It contains $a=a1$, is closed under differences and absorption, and every ideal containing $a$ contains all its multiples.

\exercise{27.14}{True and false statements}
\begin{enumerate}[label=(\alph*)]
\item False. $(0)\subset\Z$ is prime but not maximal, since $\Z$ is a domain and not a field.
\item True. The quotient by a maximal ideal is a field and hence a domain.
\item True. Every subfield of $\Q$ contains $1$, all integers, and all nonzero integer inverses, hence all rationals.
\item False. The prime subfield of $\C$ is $\Q$.
\item True. The unity generates a copy of $\Q$ in characteristic zero or of $\Z_p$ in characteristic $p$.
\item True. $\Z_2\times\Z_2$ has zero divisors and contains the diagonal field $\{(0,0),(1,1)\}$.
\item True. Send $m/n$ to $(m1)(n1)^{-1}$; characteristic zero ensures no nonzero integer maps to zero.
\item False. $(x^2)\subset F[x]$ is not prime: $x^2$ belongs but $x$ does not.
\item True. Choose a nonzero element of smallest degree and use the division algorithm; include the zero ideal separately.
\item False. $(x^2)$ is principal but its quotient has a nonzero nilpotent element $x+(x^2)$, so is not a field.
\end{enumerate}

\exercise{27.15}{A maximal ideal of an integer product}
Let $M=2\Z\times\Z$. Projection onto the first coordinate followed by reduction modulo two is onto $\Z_2$ and has kernel $M$. Thus $(\Z\times\Z)/M\cong\Z_2$, a field, and $M$ is maximal.

\exercise{27.16}{A prime ideal that is not maximal}
Let $P=\{0\}\times\Z$. Projection onto the first coordinate gives quotient $\Z$, a domain but not a field. Thus $P$ is prime and not maximal. Explicitly $P\subsetneq2\Z\times\Z\subsetneq\Z\times\Z$.

\exercise{27.17}{An ideal that is not prime}
Let $I=4\Z\times\Z$. It is proper and nonzero. Neither factor $(2,0)$ of $(2,0)(2,0)=(4,0)$ lies in $I$, but their product does. Hence $I$ is not prime.

\exercise{27.18}{A reducible rational quadratic quotient}
Since $x^2-5x+6=(x-2)(x-3)$, the defining polynomial is reducible over $\Q$. In the quotient, the two linear-factor cosets are nonzero (their degrees are too small to be multiples of the quadratic), and their product is zero. Thus the quotient is not a field.

\exercise{27.19}{An irreducible rational quadratic quotient}
The discriminant of $x^2-6x+6$ is $36-24=12$, whose square root $2\sqrt3$ is irrational. Thus the quadratic has no rational root and is irreducible over $\Q$. Its generated ideal is maximal, so the quotient is a field, isomorphic to $\Q(\sqrt3)$ by evaluating at $3+\sqrt3$.

\exercise{27.20}{Synopsis: maximality produces a field quotient}
If $M$ is maximal and $a\notin M$, the ideal $M+(a)$ must be $R$, so $1=m+ra$ for some $m\in M,r\in R$. Passing to cosets gives $(r+M)(a+M)=1+M$, so every nonzero coset has an inverse and $R/M$ is a field.

\exercise{27.21}{Synopsis: a field quotient forces maximality}
If $R/M$ is a field and $I$ is an ideal properly containing $M$, choose $a\in I\setminus M$; an inverse of $a+M$ gives $ra-1\in M$ for some $r\in R$, so $1\in I$ and $I=R$. Thus $M$ has no proper intermediate ideal and is maximal.

\exercise{27.22}{Synopsis: ideals in a polynomial ring}
For a nonzero ideal $I\subset F[x]$, choose $g\in I$ of smallest degree. Dividing any $f\in I$ by $g$ gives a remainder in $I$ of smaller degree, hence zero; therefore $I=(g)$. The zero ideal is $(0)$.

\exercise{27.23}{Synopsis: maximal polynomial ideals are irreducible}
If $(p)$ is maximal but $p=fg$ with both factors of positive degree, then $(p)\subsetneq(f)\subsetneq F[x]$, a contradiction. Hence a nonconstant generator of a maximal ideal must be irreducible.

\exercise{27.24}{Finite prime quotients}
If $P$ is prime, $R/P$ is an integral domain and is finite because $R$ is finite. For any nonzero $a$ in a finite domain, the multiplication map $x\mapsto ax$ is injective by cancellation, hence surjective. It therefore takes some $x$ to $1$, giving an inverse for $a$. Thus $R/P$ is a field, and $P$ is maximal.

\exercise{27.25}{Different cyclic-ring subrings}
It is possible when the subrings need not share the ambient unity. In $R=\Z_2\times\Z_3$, the subrings $\Z_2\times\{0\}$ and $\{0\}\times\Z_3$ are isomorphic to $\Z_2$ and $\Z_3$, respectively. Their identities are $(1,0)$ and $(0,1)$, not the ambient $(1,1)$. This distinction avoids incorrectly equating their characteristics with that of $R$.

\exercise{27.26}{Two different prime-field subrings}
The same example $\Z_p\times\Z_q$, for distinct primes $p,q$, contains the field subrings $\Z_p\times\{0\}$ and $\{0\}\times\Z_q$. Each is a field under the inherited operations, with its own identity. Thus the answer is yes under the book's subring convention.

\exercise{27.27}{The domain case}
In an integral domain, a nonzero idempotent $e$ satisfies $e(e-1)=0$, so $e=1$. Therefore the identity of any field subring is the ambient identity. A copy of $\Z_p$ forces the ambient characteristic to be $p$, and a copy of $\Z_q$ forces it to be $q$. Since $p\ne q$ is impossible, no such domain exists. Equivalently distinct nonzero orthogonal identities would create zero divisors.

\exercise{27.28}{Direct proof that maximal ideals are prime}
Suppose $ab\in M$ and $a\notin M$. The ideal $M+(a)$ properly contains $M$, so maximality makes it $R$. Hence $1=ra+m$ for some $r\in R,m\in M$. Multiply by $b$:
\[b=rab+mb\in M,\]
since both terms belong to $M$. This proves $a\in M$ or $b\in M$, and $M$ is proper by definition, so it is prime.

\exercise{27.29}{Maximal ideals and simple quotient rings}
Ideals $J$ containing $N$ correspond to ideals $J/N$ of $R/N$: inverse image under the quotient map reverses this assignment. To see the image direction, multiplication by $r+N$ sends $j+N$ to $rj+N\in J/N$; the inverse-image direction follows from the ideal test.
Therefore the only ideals of $R/N$ are zero and the whole quotient exactly when the only ideals containing $N$ are $N$ and $R$. Requiring a nontrivial quotient is equivalent to $N\ne R$. These statements say precisely that $N$ is maximal if and only if $R/N$ is simple.

\exercise{27.30}{Nonzero prime ideals of a one-variable polynomial ring}
Write the proper nonzero prime ideal as $(p)$, using the principal-ideal theorem for $F[x]$. If $p=fg$ with $\deg f,\deg g>0$, primality puts $f$ or $g$ in $(p)$, impossible by degree. Thus $p$ is irreducible. The quotient by an irreducible polynomial over a field is a field, so $(p)$ is maximal.

\exercise{27.31}{Divisibility and membership}
By definition, $(f)=\{hf:h\in F[x]\}$. Thus $g\in(f)$ means exactly that $g=hf$ for some polynomial $h$, which is the statement $f\mid g$. This includes $f=0$: membership and divisibility both then require $g=0$.

\exercise{27.32}{The ideal generated by two polynomials}
The set $N=(f,g)$ contains zero and is closed under subtraction:
\[(rf+sg)-(r'f+s'g)=(r-r')f+(s-s')g.\]
Multiplying by any $h\in F[x]$ gives $(hr)f+(hs)g$, again in $N$, so it is an ideal.
If both $f,g$ were irreducible and had distinct degrees, their only common divisors would be constants: a nonconstant common divisor of an irreducible is an associate of it, forcing equal degrees. The Euclidean algorithm then gives $uf+vg=1$, which would make $N=F[x]$. This contradicts the stated hypothesis; thus they cannot both be irreducible.

\exercise{27.33}{A one-variable Nullstellensatz}
\idea Replace the finitely generated ideal by a single greatest common divisor.
Assume the Fundamental Theorem of Algebra. Write $(f_1,\ldots,f_r)=(h)$ using the principal-ideal theorem. If $h$ is a nonzero constant, the ideal is the whole ring and the assertion is immediate. If $h=0$, all $f_i$ vanish everywhere; then $g$ vanishes everywhere, so it is the zero polynomial and belongs to the ideal.
Otherwise factor $h=c\prod_j(x-a_j)^{e_j}$. Its zeros are exactly the common zeros of the $f_i$, since $h$ divides every $f_i$ and is their polynomial linear combination. Each $a_j$ is therefore a zero of $g$. Taking $m\ge\max_j e_j$ makes $h\mid g^m$, so $g^m\in(f_1,\ldots,f_r)$.
Conversely assume the stated Nullstellensatz. If a nonconstant $f$ had no complex zero, the condition on $g=1$ would be vacuously true. It would follow that $1=1^m\in(f)$, impossible since a nonconstant polynomial cannot divide $1$. Thus every nonconstant complex polynomial has a zero.

\exercise{27.34}{The sum of two ideals}
\textbf{(a)} Zero is $0+0\in A+B$. Subtracting $a+b$ and $a'+b'$ yields $(a-a')+(b-b')\in A+B$. For $r\in R$, both $r(a+b)=ra+rb$ and $(a+b)r=ar+br$ lie in $A+B$. Thus this is an ideal even if $R$ is noncommutative.
\textbf{(b)} Since $0\in A\cap B$, write $a=a+0$ and $b=0+b$. This proves $A\subset A+B$ and $B\subset A+B$.

\exercise{27.35}{The product of two ideals}
\textbf{(a)} The zero sum can be represented as $0\cdot0$. Sums of finite sums of products remain such sums; a negative is obtained by replacing $a_i$ by $-a_i$. Also
\[r\sum_i a_ib_i=\sum_i(ra_i)b_i,\qquad(\sum_i a_ib_i)r=\sum_i a_i(b_ir),\]
and the changed factors stay in the appropriate ideals. Hence $AB$ is an ideal.
\textbf{(b)} Each $a_ib_i$ lies in $A$ by right absorption and in $B$ by left absorption. Both ideals are closed under finite sums, so $AB\subset A\cap B$.

\exercise{27.36}{The ideal quotient}
Let $I=A:B=\{r:rb\in A\text{ for every }b\in B\}$. Zero belongs to $I$. If $r,s\in I$, then $(r-s)b=rb-sb\in A$ for every $b\in B$, so $r-s\in I$. If $t\in R$, then $(tr)b=t(rb)\in A$. Thus $tr\in I$, and commutativity gives the other side too. Therefore $A:B$ is an ideal.

\exercise{27.37}{A one-sided matrix ideal}
Let $S$ be the matrices with second row zero. Differences and products of such matrices still have second row zero, so $S$ is a subring. For any $B\in M_2(F)$, the second row of $AB$ is zero whenever $A\in S$, so $S$ is a right ideal.
Using matrix units, $E_{11}\in S$ but $E_{21}E_{11}=E_{21}\notin S$. Therefore multiplication on the left by an arbitrary ambient matrix need not stay in $S$, and $S$ is not a left ideal.

\exercise{27.38}{Simplicity of a matrix ring}
Let $I$ be a nonzero two-sided ideal in $M_2(\Z_2)$, and choose $A\in I$ with nonzero entry $a_{ij}=1$. For any $r,s\in\{1,2\}$,
\[E_{ri}AE_{js}=a_{ij}E_{rs}=E_{rs}\in I.\]
Thus $I$ contains all four matrix units. Every binary matrix is the sum of the matrix units corresponding to its nonzero entries, so $I=M_2(\Z_2)$. Hence the only two-sided ideals are zero and the whole ring, proving simplicity.
